\documentclass[10pt,a4paper]{amsart}
\usepackage{amsmath,amssymb,amsthm,graphicx}
\usepackage{verbatim}
\newenvironment{reference}{\comment}{\endcomment}
%\newenvironment{reference}{}{}
\newenvironment{slogan}{\comment}{\endcomment}
\newenvironment{history}{\comment}{\endcomment}
% For commutative diagrams we use Xy-pic
\usepackage[all]{xy}

% We use 2cell for 2-commutative diagrams.
\xyoption{2cell}
% Theorem environments.
%
\theoremstyle{plain}
\newtheorem{theorem}[subsection]{Theorem}
\newtheorem{proposition}[subsection]{Proposition}
\newtheorem{lemma}[subsection]{Lemma}

\theoremstyle{definition}
\newtheorem{definition}[subsection]{Definition}
\newtheorem{example}[subsection]{Example}
\newtheorem{exercise}[subsection]{Exercise}
\newtheorem{situation}[subsection]{Situation}

\theoremstyle{remark}
\newtheorem{remark}[subsection]{Remark}
\newtheorem{remarks}[subsection]{Remarks}

\numberwithin{equation}{subsection}

% Macros
%
\def\lim{\mathop{\mathrm{lim}}\nolimits}
\def\colim{\mathop{\mathrm{colim}}\nolimits}
\def\Spec{\mathop{\mathrm{Spec}}}
\def\Hom{\mathop{\mathrm{Hom}}\nolimits}
\def\Ext{\mathop{\mathrm{Ext}}\nolimits}
\def\SheafHom{\mathop{\mathcal{H}\!\mathit{om}}\nolimits}
\def\SheafExt{\mathop{\mathcal{E}\!\mathit{xt}}\nolimits}
\def\Sch{\mathit{Sch}}
\def\Mor{\mathop{\mathrm{Mor}}\nolimits}
\def\Ob{\mathop{\mathrm{Ob}}\nolimits}
\def\Sh{\mathop{\mathit{Sh}}\nolimits}
\def\NL{\mathop{N\!L}\nolimits}
\def\CH{\mathop{\mathrm{CH}}\nolimits}
\def\proetale{{pro\text{-}\acute{e}tale}}
\def\etale{{\acute{e}tale}}
\def\QCoh{\mathit{QCoh}}
\def\Ker{\mathop{\mathrm{Ker}}}
\def\Im{\mathop{\mathrm{Im}}}
\def\Coker{\mathop{\mathrm{Coker}}}
\def\Coim{\mathop{\mathrm{Coim}}}

% Boxtimes
%
\DeclareMathSymbol{\boxtimes}{\mathbin}{AMSa}{"02}

%
% Macros for moduli stacks/spaces
%
\def\QCohstack{\mathcal{QC}\!\mathit{oh}}
\def\Cohstack{\mathcal{C}\!\mathit{oh}}
\def\Spacesstack{\mathcal{S}\!\mathit{paces}}
\def\Quotfunctor{\mathrm{Quot}}
\def\Hilbfunctor{\mathrm{Hilb}}
\def\Curvesstack{\mathcal{C}\!\mathit{urves}}
\def\Polarizedstack{\mathcal{P}\!\mathit{olarized}}
\def\Complexesstack{\mathcal{C}\!\mathit{omplexes}}
% \Pic is the operator that assigns to X its picard group, usage \Pic(X)
% \Picardstack_{X/B} denotes the Picard stack of X over B
% \Picardfunctor_{X/B} denotes the Picard functor of X over B
\def\Pic{\mathop{\mathrm{Pic}}\nolimits}
\def\Picardstack{\mathcal{P}\!\mathit{ic}}
\def\Picardfunctor{\mathrm{Pic}}
\def\Deformationcategory{\mathcal{D}\!\mathit{ef}}
\usepackage{cleveref}
\crefname{originalreference}{original reference}{original references}
\Crefname{originalreference}{Original reference}{Original references}
\makeatletter
\newlabel{algebra-lemma-bound-depth}{{algebra-lemma-bound-depth}{source}{}{sourcepointer-1}{}}
\newlabel{algebra-lemma-bound-depth@cref}{{[originalreference][0][]algebra-lemma-bound-depth}{[1][0][]source}{}{}{}}
\makeatother
\crefname{originallemma}{original lemma}{original lemmas}
\Crefname{originallemma}{Original lemma}{Original lemmas}
\makeatletter
\newlabel{lemma-dualizing}{{lemma-dualizing}{source}{}{sourcepointer-2}{}}
\newlabel{lemma-dualizing@cref}{{[originallemma][0][]lemma-dualizing}{[1][0][]source}{}{}{}}
\makeatother
\crefname{originallemma}{original lemma}{original lemmas}
\Crefname{originallemma}{Original lemma}{Original lemmas}
\makeatletter
\newlabel{lemma-sitting-in-degrees}{{lemma-sitting-in-degrees}{source}{}{sourcepointer-3}{}}
\newlabel{lemma-sitting-in-degrees@cref}{{[originallemma][0][]lemma-sitting-in-degrees}{[1][0][]source}{}{}{}}
\makeatother
\title{Cohen-Macaulayness read off from the vanishing of Ext into the normalized dualizing complex}
\author{The Stacks Project Authors}
\date{stacks-project:dualizing.tex:local-snapshot-20261012; GNU FDL 1.2 or later}
\begin{document}
\maketitle
\noindent\textit{Editorial note.} All delivered mathematical text is the author's own, copied from the pinned Stacks Project chapter source sources/stacks-dualizing/pinned-original-source.tex; no statement or proof step was rewritten, reordered or removed.\par\medskip
\noindent\textit{Editorial note.} Preamble provenance: the Stacks Project keeps one single shared preamble for every chapter; this card takes it directly from the staged sources/stacks-dualizing/preamble.tex as three verbatim spans (lines 8-12, 14-18 and 154-217), so no derived file was written and nothing was added or redefined. The wrapper class and the base packages are supplied by the assembler.\par\medskip
\noindent\textit{Editorial note.} Bibliography provenance: neither the counted statement nor its proof contains a citation (this lane rescans the whole span before writing the card), so the card carries no bibliography and no entry was reconstructed.\par\medskip
\noindent\textit{Editorial note.} Selection provenance: this card was opened under the file-based control contract (CONTROL.md), with the mechanical pre-card gate run on its statement and proof interval before the card was written.\par\medskip
\noindent\textit{Allowed standard prerequisites (editorial).}\begin{itemize}
\item The degrees in which Ext into the normalized dualizing complex can be nonzero. Invoked at the start of the counted proof (source lines 3323-3324). Confines the Ext modules to a single relevant degree and fixes the comparison with the depth.
\item The double dual statement identifying a finite module with the internal hom of its Ext into the dualizing complex. Invoked inside the counted proof (source lines 3331-3335). Turns a single nonzero Ext into an identification of the module as an Ext of the resulting module.
\item The bound of the depth of a finite module by the dimension of its support. Invoked inside the counted proof (source lines 3338-3339). Supplies the reverse inequality that upgrades the single-nonzero condition to Cohen-Macaulayness.
\end{itemize}\medskip
\ifdefined\hypertarget\hypertarget{sourcepointer-1}{}\fi
\noindent\textit{Source reference.} Original reference algebra-lemma-bound-depth: sources/stacks-dualizing/pinned-original-source.tex, source lines 3339--3339. Navigation pointer only: the counted text refers to the Stacks item algebra-lemma-bound-depth, which belongs to another chapter of the project and is not part of this card. The pointer carries the target's own label as its number because this chapter's staged snapshot does not contain the other chapter's tag table; its content is not reproduced here.\par\medskip
\ifdefined\hypertarget\hypertarget{sourcepointer-2}{}\fi
\noindent\textit{Source reference.} Original lemma lemma-dualizing: sources/stacks-dualizing/pinned-original-source.tex, source lines 2818--2818. Navigation pointer only: the counted text refers to the Stacks item lemma-dualizing, which is an existing labeled result of this same chapter and is not part of this card. The pointer carries the target's own label as its number and its definition line as its span; its content is not reproduced here.\par\medskip
\ifdefined\hypertarget\hypertarget{sourcepointer-3}{}\fi
\noindent\textit{Source reference.} Original lemma lemma-sitting-in-degrees: sources/stacks-dualizing/pinned-original-source.tex, source lines 3232--3232. Navigation pointer only: the counted text refers to the Stacks item lemma-sitting-in-degrees, which is an existing labeled result of this same chapter and is not part of this card. The pointer carries the target's own label as its number and its definition line as its span; its content is not reproduced here.\par\medskip
\begin{lemma}
\label{lemma-local-CM}
Let $(A, \mathfrak m, \kappa)$ be a Noetherian local ring
with normalized dualizing complex $\omega_A^\bullet$. Let $M$
be a finite $A$-module. The following are equivalent
\begin{enumerate}
\item $M$ is Cohen-Macaulay,
\item $\Ext^i_A(M, \omega_A^\bullet)$ is nonzero for at most one $i$,
\item $\Ext^{-i}_A(M, \omega_A^\bullet)$ is zero for
$i \not = \dim(\text{Supp}(M))$.
\end{enumerate}
Denote $CM_d$ the category of finite Cohen-Macaulay $A$-modules
of depth $d$. Then $M \mapsto \Ext^{-d}_A(M, \omega_A^\bullet)$
defines an anti-auto-equivalence of $CM_d$.
\end{lemma}

\begin{proof}
We will use the results of Lemma \ref{lemma-sitting-in-degrees}
without further mention. Fix a finite module $M$.
If $M$ is Cohen-Macaulay, then only
$\Ext^{-d}_A(M, \omega_A^\bullet)$ can be nonzero,
hence (1) $\Rightarrow$ (3).
The implication (3) $\Rightarrow$ (2) is immediate.
Assume (2) and let $N = \Ext^{-\delta}_A(M, \omega_A^\bullet)$
be the nonzero $\Ext$ where $\delta = \text{depth}(M)$. Then, since
$$
M[0] = R\Hom_A(R\Hom_A(M, \omega_A^\bullet), \omega_A^\bullet) =
R\Hom_A(N[\delta], \omega_A^\bullet)
$$
(Lemma \ref{lemma-dualizing})
we conclude that $M = \Ext_A^{-\delta}(N, \omega_A^\bullet)$.
Thus $\delta \geq \dim(\text{Supp}(M))$. However,
since we also know that $\delta \leq \dim(\text{Supp}(M))$
(Algebra, Lemma \ref{algebra-lemma-bound-depth}) we conclude that $M$ is
Cohen-Macaulay.

\medskip\noindent
To prove the final statement, it suffices to show that
$N = \Ext^{-d}_A(M, \omega_A^\bullet)$ is in $CM_d$
for $M$ in $CM_d$. Above we have seen that
$M[0] = R\Hom_A(N[d], \omega_A^\bullet)$ and this proves the
desired result by the equivalence of (1) and (3).
\end{proof}

\end{document}
